\begin{appendices}

\section{支撑材料清单}

\begin{table}[!htbp]
  \caption{支撑材料目录说明}\centering
  \begin{tabularx}{\textwidth}{>{\raggedright\arraybackslash}p{4.6cm} L}
    \toprule[1.5pt]
    文件／目录 & 内容 \\
    \midrule[1pt]
    \texttt{code/model.py} & 有界误差几何：楔形半平面、多边形裁剪、半平面交的空／无界／有界判定、直径、最小包围圆、光学格、测站布局生成与覆盖校验 \\
    \texttt{code/optimized.py} & 在线策略主体：轨迹维护、共享测站扫描、认证清除与光学兜底入口 \\
    \texttt{code/q3\_policy.py} & 第三问专用：目标圆外切约束、\texttt{no\_signal} 圆盘排除、第二测点、圆心补测、混合贪心调度 \\
    \texttt{code/q4\_policy.py} & 第四问专用：$22$ 站列表校验与凸包方向覆盖检查、选择性重测、有界等待、补测精化、条带清除 \\
    \texttt{code/q4\_optical.py} & 认证条带覆盖与扫描路线成本比较 \\
    \texttt{code/\allowbreak practice\_robot.py} & 官方模拟器 HTTP 客户端：串行发送、超时原样重试、未决动作保护、逐动作日志与预算控制 \\
    \texttt{code/config.json} & 当前全部算法参数与 $22$ 站坐标 \\
    \texttt{code/layout\_opt.py}、\allowbreak\texttt{code/certify22.py} & 方向覆盖余量计算与 $h=1\mmet$ 网格证书 \\
    \texttt{official\_practice/} & 演练样本：\texttt{q3\_runs.csv}、\texttt{q4\_runs.csv} 为 $50+70$ 局的逐局明细，
      \texttt{official\_runs.json} 另含被排除记录的清单与原因 \\
    \texttt{logs/}（三次正式测试） & 模拟器导出的加密行为日志，文件名未作任何修改 \\
    \texttt{AI工具使用详情.pdf} & 按 2026 年 AI 工具使用规定提交的详情说明 \\
    \bottomrule[1.5pt]
  \end{tabularx}
\end{table}

\noindent 运行方式：安装 Python 3.10 及以上与 \texttt{numpy}、\texttt{scipy}、\texttt{requests}，在模拟器接口就绪后执行
\begin{center}
\texttt{python practice\_robot.py --question 3 --robot-id <队号> --case-code <案例编码> --connect}
\end{center}
（第四问将 \texttt{--question} 改为 \texttt{4}）。不加 \texttt{--connect} 时程序不会发送任何请求，便于离线检查。

\section{关键参数配置}

\begin{table}[!htbp]
  \caption{算法参数及其来源}\centering\small
  \begin{tabularx}{\textwidth}{@{}l l L l@{}}
    \toprule[1.5pt]
    参数 & 取值 & 作用 & 来源 \\
    \midrule[1pt]
    \texttt{region\_radius\_m}        & $1800$   & 目标区域半径 & 题目 \\
    \texttt{receiver\_radius\_min\_m} & $1000$   & 接收半径下界，覆盖证明的基准 & 题目 \\
    \texttt{receiver\_radius\_max\_m} & $1500$   & 接收半径上界，初始楔形长度 & 题目 \\
    \texttt{bearing\_bound\_deg}      & $1.005$  & 保守误差界 & \cref{asu:bounded} \\
    \texttt{clear\_m}                 & $20$     & 认证清除半径 & 题目 \\
    \texttt{near\_m}                  & $5$      & 近场阈值 & 题目 \\
    \texttt{optical\_grid\_m}         & $25$     & 光学方格边长，$25/\sqrt2<20$ & \cref{eq:grid} \\
    \texttt{q3\_survey\_layout}       & hexagon  & 第三问七站布局 & \cref{thm:seven} \\
    \texttt{q3\_ring\_radius\_m}      & $1150$   & 第三问环半径，可行区间 $[1122.96,1732.05]$ & \cref{cor:ring} \\
    \texttt{q3\_policy}               & mixed    & 搜索与目标混合贪心调度 & \cref{sec:q3:shared} \\
    \texttt{q3\_second\_forward/lateral\_m} & $200$／$100$ & 第二测点偏移 & \cref{sec:q2} \\
    \texttt{q3\_refine\_max\_m}       & $300$    & 圆心补测的适用上限 & \cref{sec:q3:refine} \\
    \texttt{q3\_refine\_offset\_m}    & $60$     & 补测横移距离 & \cref{sec:q3:refine} \\
    \texttt{q4\_policy}               & mixed\_ring & 第四问联合调度 & \cref{sec:q4:online} \\
    \texttt{q4\_station\_list}        & $22$ 点  & 认证布局，删除即回退 $25$ 站 & \cref{tab:q4layout} \\
    \texttt{q4\_selective\_remeasure} & true     & 选择性重测 & \cref{eq:selective} \\
    \texttt{q4\_remeasure\_min\_sine} & $0.3$    & 重测的最小交会角正弦 & \cref{eq:selective} \\
    \texttt{q4\_remeasure\_min\_gap\_m}& $300$   & 重测的最小测点间距 & \cref{eq:selective} \\
    \texttt{q4\_station\_help\_m}     & $950$    & 判定测站对目标是否有价值 & \cref{eq:selective} \\
    \texttt{q4\_wait\_for\_survey}／\texttt{q4\_wait\_mode} & true／inner & 有界等待，仅限内环发现的目标 & \cref{sec:q4:online} \\
    \texttt{q4\_wait\_stale\_m}       & $1400$   & 超过该距离不再等待 & \cref{sec:q4:online} \\
    \texttt{q4\_refine\_max\_m}       & $150$    & 第四问补测的适用上限 & \cref{sec:q4:online} \\
    \texttt{q4\_strip\_cover}         & true     & 认证条带覆盖 & \cref{eq:strip} \\
    \texttt{q4\_scan\_cost\_gate}     & true     & 扫描成本门控 & \cref{eq:routecost} \\
    \texttt{q4\_immediate\_near}      & true     & 收到 \texttt{near} 立即清除 & \cref{sec:q4:online} \\
    \bottomrule[1.5pt]
  \end{tabularx}\\[2pt]
  \footnotesize \texttt{config.json} 中另有若干历史参数（如通用的 \texttt{second\_forward\_m}$=650$、\texttt{second\_lateral\_m}$=250$
  以及 \texttt{q3\_grid\_m}、\texttt{q4\_grid\_m}、\texttt{q4\_triangle\_spacing\_m}）。它们\textbf{不被当前默认路径读取}：
  第三、四问的第二测点分别取 \texttt{q3\_}、\texttt{q4\_} 前缀的同名参数。保留这些条目是为了能一键回退到历史对照版本
  （例如置 \texttt{q3\_policy=legacy} 或删除 \texttt{q4\_station\_list}），复现旧轮次实验时仍需要它们，因此没有删除。
\end{table}

\section{调试记录：一次边界顶点的浮点失稳}\label{app:debug}

在第三问圆盘排除约束的首轮验证中，种子 $2031091108$ 的案例报出
\texttt{Localization certificate excluded true source by 9.7374 m}：该局虽然最终清除了全部 $15$ 个源，但中间的几何证书已经失效。根因是重复施加圆盘排除时，一个理论上位于边界的交点被浮点判为落在圆盘内部，相应二次方程的根又略微超出线段参数范围，导致边界顶点被错误丢弃，多次裁剪后可行域出现不保守的收缩。修复后改用平方距离容差保留边界顶点、对接近端点的根做范围容许与夹取、点数不足或凸包退化时保留输入多边形，并补充了重复约束的回归测试。本例说明“结果全部清除”不能替代对中间几何不变量的检查（正文\cref{sec:q3:loc}）。

\section{虚拟时间上界的推导}\label{app:bound}

所有行动点都在以原点为心、$2000\mmet$ 为半径的圆内，任意两点距离不超过 $4000\mmet$。$7$ 个测站每站至多 $20$ 次含切频的检测（$120\msec$）；至多 $16$ 个目标，每个目标至多两段全局跳转、一条保守长度 $6500\mmet$ 的光学路线、若干次检测与至多 $183$ 次清除尝试\footnote{$2000\mmet$：最外环测站半径 $1150\mmet$（第三问）与 $1866\mmet$（第四问）均不超过 $1866\mmet$，而所有源都在 $R=1800\mmet$ 内，第二测点与光学格点也仅在其附近，故所有行动点到原点的距离不超过 $2000\mmet$。$6500\mmet$：对任意“交会后凸子集”，按 $25\mmet$ 方格的列式蛇形路线，每列内部移动不超过 $50\mmet$，相邻非空列间距不超过 $\sqrt{25^2+50^2}\mmet$，且连续凸域不会跳过内部列，故单目标光学路径以 $6500\mmet$ 为上界。$183$：在 $\bar\varepsilon=1.005^\circ$、格距 $25\mmet$ 下，单个可行域的候选格中心至多 $61$ 个横向列 $\times3$ 个纵向行 $=183$ 格。三者均按最坏情形取保守上界，只用于确认 $T$ 远小于虚拟时间上限，不参与任何决策。}。于是
\begin{equation}
  T\le 7\Bigl(\frac{4000}{5}+120\Bigr)+16\Bigl[\frac{2\times4000+6500}{5}+6\times6+183\times3+5\Bigr]
  \approx6.2\times10^{4}\msec\ \ll\ 3.6\times10^{5}\msec,
\end{equation}
远小于虚拟时间上限，也说明算法不可能因死循环而被强制终止。实际演练中整局虚拟时间为 $3301\pm291\msec$，第三问的现实程序运行时间仅 $0.4\sim2.2\msec$，相对 $1200\msec$ 的限制有三个数量级的余量。


\section{被否决改进路线的完整记录}\label{app:ablation}

下列两表汇总了第三、四问中完成了完整配对实验、但最终未进入正式方案的全部改进路线。每一行都使用与开发阶段不重叠的留出案例、同一误差场做配对比较，清除率均为 $100\%$。保留这些记录的目的有两个：它们构成本文的消融实验，说明最终算法的每一项选择都有对照；同时也提醒读者，平均值改善并不足以作为采用一项改动的充分理由。

\begin{table}[!htbp]
  \caption{第三问被否决的改进方案（均完成完整配对实验，清除率均为 $100\%$）}\label{tab:q3reject}\centering
  \begin{tabularx}{\textwidth}{>{\raggedright\arraybackslash}p{4.4cm} c c L}
    \toprule[1.5pt]
    方案 & 对照 s/源 & 候选 s/源 & 否决原因 \\
    \midrule[1pt]
    两步前瞻调度：选动作时同时计入下一步的最优代价，$J(a)=C(p,a)+\lambda\min_bC(a,b)$ & $263.920$ & $292.795$ & 总移动距离由 $12583$ 升至 $14452\mmet$，$\mathrm{P}_{95}$ 同步恶化；复杂度上升而性能下降 \\
    非凸可行域：把 \texttt{no\_signal} 排除掉的圆盘真的“挖掉”，维护非凸单元集合 & $263.920$ & $264.943$ & 单元碎片化，平均光学格数由 $100.1$ 升至 $102.3$，区域缩小没有转化为更短路径 \\
    自适应第二测点：在 $5\times4$ 个离散偏移中按收益选点 & $263.920$ & $265.518$ & 均值与最坏值都变差，详见\cref{sec:q2} \\
    主动搜索：按“新增覆盖收益除以移动代价”自行选择下一测站 & $251.70$ & $260.95$ & 提前减少搜索后，已发现目标拿到的方位不足，光学成本反而上升 \\
    \bottomrule[1.5pt]
  \end{tabularx}
\end{table}

\begin{table}[!htbp]
  \caption{第四问被否决的改进方案（均完成完整实验，清除率都是 $100\%$）}\label{tab:q4reject}\centering\small
  \begin{tabularx}{\textwidth}{>{\raggedright\arraybackslash}p{3.6cm} L}
    \toprule[1.5pt]
    方案 & 实验结果与否决原因 \\
    \midrule[1pt]
    绕行代价阈值调度 &
      只有当“插入该目标所多走的路” $d(x,g)+d(g,s)-d(x,s)$ 低于阈值时才立即处理，否则推迟到收尾阶段。
      在 $27$ 站与 $25$ 站上各做 $1050$ 次仿真，最好的阈值使平均时间恶化 $8.54\%$ 与 $9.15\%$，$100$ 个随机案例中 $89$ 个退步。
      返回主路线的移动确实降了 $16\%$，但光学扫描移动增加 $49.7\%$、目标测量移动增加 $88.0\%$。 \\
    按价值门控裁剪重复测量 &
      测量次数减少 $14.36\%$、平均时间改善 $3.35\%$，但光学扫描距离增加 $39\%$，最坏案例多花 $296.49\msec$。
      裁剪过度会让少数困难源失去关键方位而被迫进入长距离搜索。 \\
    为上一条加装风险保护 &
      先检查预计光学路径长度与尝试次数，超限则强制测量。最佳参数把光学距离增幅压到 $2.17\%$、平均改善 $1.72\%$，
      却出现新的最坏案例（多花 $950.86\msec$）。\textbf{平均值改善抵不过尾部风险，整条路线放弃。} \\
    固定巡回路线与两交换改进 &
      当前“每次去最近未访问站”的路线长 $18233.1\mmet$；最近邻加两交换、$100$ 个多起点加两交换给出的都是同样的数值，
      现有路线已与这些方法找到的最好解持平。 \\
    低价值测站逐一删除 &
      在 $1120$ 个案例上统计每站的发现贡献与到达成本，再逐站做留一实验。虽有若干外环站平均价值偏低，
      但每个外环站都能构造出独占方向覆盖的反例，\textbf{没有任何站拿得到连续域的安全删除证书}。 \\
    运行时动态跳站 &
      只在“剩余测站仍能对所有未发现频道提供完整覆盖证明”时才允许跳过。$1000$ 个案例中安全证书授权的跳站次数为 $\bm 0$：
      证书足够安全但过于保守，没有实际收益。 \\
    主动搜索 &
      用保守单元维护位置--方向联合证书，按覆盖收益自行选站。原型曾有“借助真实源数决定补测”与“用有限探针代替连续证明”
      两处研究有效性缺陷，修正后仍未胜过现行调度。 \\
    \bottomrule[1.5pt]
  \end{tabularx}
\end{table}

\section{灵敏度分析的两张数据表}\label{app:sens}

本附录汇总正文中因篇幅而移出的数据表：\cref{tab:q3six}为第三问六站布局的实验记录，\cref{tab:q4search}为第四问布局搜索考察过的候选，另两张表是\cref{sec:eval}灵敏度分析所依据的原始数据。

\begin{table}[!htbp]
  \caption{六站布局实验（独立入口，默认方案未采用）}\label{tab:q3six}\centering\small
  \begin{tabularx}{\textwidth}{@{}L c c c c c@{}}
    \toprule[1.5pt]
    布局 & 最坏未覆盖距离/m & 单源漏检概率 & $120$ 例均值 s/源 & 与 $7$ 站配对差 & 更快局数 \\
    \midrule[1pt]
    $7$ 站（默认） & $988.5$（有保证） & $0$ & $258.6$ & --- & --- \\
    $6$ 站环 $960\mmet$   & $1102.0$ & $\approx4\times10^{-4}$ & $256.1$ & $-2.5$ & $69/120$（漏 $1$ 源） \\
    $6$ 站环 $1000\mmet$  & $1062.9$ & $\approx2\times10^{-4}$ & $257.2$ & $-1.4$ & $61/120$（无漏） \\
    $6$ 站环 $1040\mmet$  & $1043.5$ & $\approx8\times10^{-5}$ & $257.9$ & $-0.7$ & $63/120$（无漏） \\
    \bottomrule[1.5pt]
  \end{tabularx}
\end{table}

\begin{table}[!htbp]
  \caption{布局搜索中考察过的候选（严格证书取 $h=1\mmet$、$\delta=0.7071\mmet$；余量为负表示存在无法覆盖的点）}\label{tab:q4search}\centering\small
  \begin{tabularx}{\textwidth}{@{}L c Y L@{}}
    \toprule[1.5pt]
    布局 & 站数 & 粗网格最小余量/m & 严格证书 \\
    \midrule[1pt]
    圆心 $+12@999+12@1863.5$（上一版） & $25$ & $0.0$（解析三角剖分证明成立） & 通过 \\
    \textbf{圆心 $+9@990+12@1866$（采用）} & $\bm{22}$ & $\bm{2.42}$ & \textbf{通过，最小余量 $1.73$} \\
    圆心 $+8@990+12@1866$ & $21$ & $-0.5$ & 不通过 \\
    圆心 $+7@990+12@1866$ & $20$ & $-29$ & 不通过 \\
    三核 $+6@1050+12@1866$ 等 & $20\sim22$ & $-8\sim-43$ & 不通过 \\
    随机局部搜索所得布局 & $19\sim22$ & $-42\sim-92$ & 不通过 \\
    圆心 $+6@1200+8@2000$ 等稀疏布局 & $14\sim19$ & $-176\sim-596$ & 不通过 \\
    \bottomrule[1.5pt]
  \end{tabularx}\\[2pt]
  \footnotesize 本表是一次有限搜索的记录，覆盖“圆心 $+$ 正内环 $+$ 正外环”族及若干由该族出发的局部扰动，
  \textbf{不构成对任意 $21$ 站布局的穷举}。搜索脚本与逐布局余量见支撑材料 \texttt{layout\_opt.py}。
\end{table}


\begin{table}[!htbp]
  \caption{测向误差幅度的灵敏度（同一批种子，单位：虚拟秒/源，均 $100\%$ 清除）}\label{tab:senserr}\centering
  \begin{tabular}{cccc}
    \toprule[1.5pt]
    实际误差幅度 & 算法误差界 & 问题三 & 问题四 \\
    \midrule[1pt]
    $\pm0.5^\circ$ & $0.505^\circ$ & $732.38$ & $966.58$ \\
    $\pm1.0^\circ$ & $1.005^\circ$ & $741.00$ & $987.94$ \\
    $\pm1.5^\circ$ & $1.505^\circ$ & $745.92$ & $1015.95$ \\
    \bottomrule[1.5pt]
  \end{tabular}
\end{table}

\begin{table}[!htbp]
  \caption{方向覆盖余量对最小接收半径的灵敏度（$20\mmet$ 粗网格；$-\infty$ 表示存在无法覆盖的点）}\label{tab:sensrho}\centering
  \begin{tabular}{lccccccc}
    \toprule[1.5pt]
    $\rho_{\min}/\mmet$ & $1000$ & $998$ & $996$ & $994$ & $990$ & $985$ & $980$ \\
    \midrule[1pt]
    $22$ 站（采用）余量$/\mmet$ & $2.417$ & $2.417$ & $2.417$ & $0.482$ & $-\infty$ & $-\infty$ & $-\infty$ \\
    $25$ 站（解析）余量$/\mmet$ & $0.000$ & $-\infty$ & $-\infty$ & $-\infty$ & $-\infty$ & $-\infty$ & $-\infty$ \\
    \bottomrule[1.5pt]
  \end{tabular}
\end{table}

\begin{table}[!htbp]
  \caption{第四问历次结构性改进的配对实验（单位：虚拟秒/源，两版均 $100\%$ 清除）}\label{tab:q4paired}\centering
  \begin{tabular}{lcccc}
    \toprule[1.5pt]
    改进 & 数据集 & 对照 & 候选 & 下降 \\
    \midrule[1pt]
    $49$ 站方格 $\to$ $27$ 站三角网格 & $100$ 例，$1310$ 源 & $928.46$ & $767.64$ & $17.32\%$ \\
    （同上，压力集） & $50$ 例，$670$ 源 & $916.49$ & $727.66$ & $20.60\%$ \\
    $27$ 站 $\to$ $25$ 站双环 $+$ 混合调度 & $100$ 例随机 & $792.16$ & $566.30$ & $28.51\%$ \\
    选择性重测 $+$ 有界等待 $+$ 精化 & $220$ 例，$2908$ 源 & $534.05$ & $499.98$ & $6.38\%$ \\
    认证条带 $+$ \texttt{near} 即清 $+$ 成本门控 & $450$ 例，$5935$ 源 & $499.45$ & $493.55$ & $1.18\%$ \\
    \bottomrule[1.5pt]
  \end{tabular}
\end{table}

\begin{table}[!htbp]
  \caption{第三问历次结构性改进的配对实验（单位：虚拟秒/源，两版均 $100\%$ 清除）}\label{tab:q3paired}\centering
  \begin{tabular}{lcccc}
    \toprule[1.5pt]
    改进 & 数据集 & 对照 & 候选 & 下降 \\
    \midrule[1pt]
    单次测向 $\to$ 两次测向 & $40$ 例，$498$ 源 & $977.71$ & $764.71$ & $21.79\%$ \\
    动态处理顺序 $+$ 共享测站 & $100$ 例，$1310$ 源 & $723.82$ & $653.11$ & $9.77\%$ \\
    $25$ 站网格 $\to$ $7$ 站解析覆盖 & $100$ 例，$1308$ 源 & $654.53$ & $346.62$ & $47.04\%$ \\
    环半径 $1150$ $+$ 混合贪心调度 & $100$ 例，$1315$ 源 & $343.73$ & $263.86$ & $23.24\%$ \\
    （同上，压力集） & $70$ 例，$890$ 源 & $336.31$ & $281.83$ & $16.20\%$ \\
    圆心补测精化 & $100$ 例留出 & $262.6$ & $260.0$ & $1.0\%$ \\
    \bottomrule[1.5pt]
  \end{tabular}\\[2pt]
  \footnotesize 各行使用的案例集合不同，\textbf{不能把各行降幅连乘}得到一个“总提升”。
\end{table}

\section{核心源程序}

以下为完整的可运行源程序。为便于核对，文件内容与支撑材料中的副本逐字节一致。

\subsection{有界误差几何与测站布局：\texttt{model.py}}
\lstinputlisting[language=Python]{code/model.py}

\subsection{在线策略主体：\texttt{optimized.py}}
\lstinputlisting[language=Python]{code/optimized.py}

\subsection{第三问专用策略：\texttt{q3\_policy.py}}
\lstinputlisting[language=Python]{code/q3_policy.py}

\subsection{第四问专用策略：\texttt{q4\_policy.py}}
\lstinputlisting[language=Python]{code/q4_policy.py}

\subsection{认证条带覆盖：\texttt{q4\_optical.py}}
\lstinputlisting[language=Python]{code/q4_optical.py}

\subsection{方向覆盖余量与网格证书：\texttt{layout\_opt.py}、\texttt{certify22.py}}
\lstinputlisting[language=Python]{code/layout_opt.py}
\lstinputlisting[language=Python]{code/certify22.py}

\subsection{模拟器通信客户端：\texttt{practice\_robot.py}}
\lstinputlisting[language=Python]{code/practice_robot.py}

\subsection{参数配置：\texttt{config.json}}
\lstinputlisting{code/config.json}

\end{appendices}
